\begingroup
\section{The original relation graph supplies the simultaneous mixed flags}
\label{sec:NFI}

This calculation joins the original-data relation formulas NTG4--22 to
the simultaneous restriction and quotient formulas STF4--22. The finite
Hermitian operations are proved below; the Schur quotient operations
being transported are those of Boyd and Vandenberghe,
Appendix A.5.5, p.~650, and C.4.1, p.~673
\cite{human:boyd-vandenberghe}, with their complex specialization
proved in STF8--12. Every scalar polynomial and norm retains the
Meixner--Pollaczek parameter and mass dictionary of NTG3
\cite[eqs.~(18.19.8)--(18.19.9), (18.22.8)]{human:dlmf18}.

Use NTG's original $A$, $x_r$ and $\Omega_R$, and choose an integer
$0\le R<q$. The symbol $D^{\rm nat}_n$ in this section denotes the
matrix called $D_n$ in NTG18. STF2's target-cutoff convention is
$D_N=D^{\rm nat}_{N-v}$; therefore $N=q+r$ corresponds to $n=L+r$.
Define, with the unique positive Hermitian square roots,
\[
 G=\overline{\Omega_R},\qquad H_1=A^*GA=D_{q+R},\qquad
 H=\operatorname{diag}(H_1,4),\qquad
 U=G^{1/2}AH_1^{-1/2},\quad U_4=\operatorname{diag}(U,4).
                                                               \tag{NFI1}
\]
All matrices in the inverses are positive definite by NTG12,16,20.
Direct multiplication gives $U^*U=I_m$. It is square and injective,
so its inverse is $U^*$ and $UU^*=I_m$ as well. In particular
$U_4$ is a unitary map between the two stated Euclidean coordinate
spaces, not an alteration of the original source measure.

For $R\le t\le q$ set $d=t-R$ and retain
\[
 Y_t=[x_{R+1},\ldots,x_t],\qquad
 Z_t=G^{-1/2}\overline{Y_t},\qquad
 B_t=Z_t^*U=Y_t^{\mathsf T}AH_1^{-1/2},\qquad
 \Theta_t=\operatorname{diag}_{r=R+1}^{t}(-i^{-(L+r)}).
                                                               \tag{NFI2}
\]
Here $Y_t,Z_t$ are $m$ by $d$, $B_t$ is $d$ by $m$, and
$\Theta_t$ is unitary of size $d$. For $d=0$ each displayed
rectangular matrix is the corresponding empty map. The equality
for $B_t$ follows from $(G^{-1/2})^*=G^{-1/2}$ and the formula
for $U$; it keeps the ordinary transpose required by the original
complex-linear functionals.

Extend STF4's endpoint notation to every $R\le t\le q$ by
$C_t=[U_{R+1}^{\mathsf T},\ldots,U_t^{\mathsf T}]^{\mathsf T}H^{-1/2}$,
where the bracket means the vertical stack of the original row-block
matrices $U_r$, and use the empty stack for $t=R$.
It agrees with STF4 at its two stated endpoints.
These four-copy original row updates are stacked by degree. Write their
row coordinates as $(r,b)$, with $R+1\le r\le t$ and $1\le b\le4$,
and let $\Pi_t$ be the permutation unitary sending coordinate
$(r,b)$ to $(b,r)$. Thus it changes only the stated row ordering.
Set
\[
 \mathscr R_t=\operatorname{diag}(\Theta_t^*,4)\Pi_t,\qquad
 \mathscr Z_t=\operatorname{diag}(Z_t,4).
                                                               \tag{NFI3}
\]
NTG22 gives the full original row
$\rho_{L+r}/\sqrt{1+\zeta_{L+r}}=-i^{-(L+r)}x_r^{\mathsf T}A$.
Multiplying on the right by $H_1^{-1/2}$, stacking and using the
defined permutation proves, entry by entry,
\[
 \Pi_t C_t=\operatorname{diag}(\Theta_tB_t,4),\qquad
 \boxed{\mathscr R_t C_t=\mathscr Z_t^*U_4.}
                                                               \tag{NFI4}
\]
The inverse row map is $\mathscr R_t^*$. Thus the original native
phase has been recorded in an explicit invertible map before taking
any Gram matrix. No phase is inferred merely from equality of norms.

Taking the Gram of NFI4 gives the exact coefficient-space relation
\[
 \mathscr N_t=\mathscr Z_t\mathscr Z_t^*
       =\operatorname{diag}(G^{-1/2}\overline{Y_t}Y_t^{\mathsf T}
                                      G^{-1/2},4),\qquad
 M_t=U_4^*\mathscr N_tU_4.
                                                               \tag{NFI5}
\]
This is an equality in the original STF coordinates after the
specified unitary map. For example,
\[
 \det(I_{4m}+M_t)
       =\left(\frac{\det\Omega_t}{\det\Omega_R}\right)^4.
                                                               \tag{NFI6}
\]
To prove NFI6, $\Omega_t=\Omega_R+Y_tY_t^*$, so conjugating
by $\Omega_R^{-1/2}$ gives
$\det(I+\Omega_R^{-1/2}Y_tY_t^*\Omega_R^{-1/2})
=\det\Omega_t/\det\Omega_R$.
Complex conjugation gives the determinant of the single block of
$I+\mathscr N_t$; both determinants are positive real numbers.
There are four blocks. This also proves the empty-tail case, with
ratio $1$. It is consistent with both allocations in STF17--18,
without selecting either allocation.

For every original coefficient subspace $X$, including those in the
labelled eight-flag tuple STF16, retain its fixed independent frame $X_0$
and put
\[
 \widehat Z_X=
  \operatorname{diag}(G^{1/2}A,4)X_0(X_0^*HX_0)^{-1/2},
 \qquad \widehat P_X=\widehat Z_X\widehat Z_X^*.
                                                               \tag{NFI7}
\]
If $X=0$ use the empty frame. NFI1 and STF5 give
$\widehat Z_X=U_4Z_X$, hence
$\widehat Z_X^*\widehat Z_X=I$ and
$\widehat P_X=U_4P_XU_4^*$. This supplies the exact transported
flag with its full original conductor-dependent coefficient matrix
$A$. Its inverse is $U_4^*$ on the transformed Euclidean space,
followed by $H^{-1/2}$ when recovering original coefficient vectors.
Consequently NFI4 proves
\[
 \mathscr R_tK_{X,t}\mathscr R_t^*
       =\mathscr Z_t^*\widehat P_X\mathscr Z_t,
 \qquad
 F_t(X)=\log\det(I+\mathscr Z_t^*\widehat P_X\mathscr Z_t).
                                                               \tag{NFI8}
\]
The determinant equality uses a unitary conjugation of the entire
row space, including its zero eigenvalues. Its version on the fixed
$4m$-dimensional coefficient space is
\[
 \widehat K_{X,t}=\mathscr N_t^{1/2}\widehat P_X
                                  \mathscr N_t^{1/2}.
                                                               \tag{NFI9}
\]
For completeness define the map
$\mathcal W_t(\mathscr N_t^{1/2}z)=\mathscr Z_t^*z$ on
$\operatorname{im}\mathscr N_t^{1/2}$, extending it by zero on
the kernel. It is well defined and isometric because
$\|\mathscr Z_t^*z\|^2=z^*\mathscr N_tz
=\|\mathscr N_t^{1/2}z\|^2$.
Then the right side of NFI8 equals
$\mathcal W_t\widehat K_{X,t}\mathcal W_t^*$.
The inverse on the image is $\mathcal W_t^*$; its complementary
kernel contributes only zero eigenvalues. This proves the precise
spectral correspondence used by both logarithmic functions in STF,
including rank loss and coincident flags.

There is a corresponding formula for every original quotient
$K\subset Y$. Retain $Q_0$ and $E$ from STF8 and set
$\widehat E=U_4E$; then
$\widehat E\widehat E^*=\widehat P_Y-\widehat P_K$.
Substituting NFI4 into the complete attained minimum STF10 gives
\[
 Q_t=Q_0^{1/2}\left[I+
 \widehat E^*\mathscr Z_t
 (I+\mathscr Z_t^*\widehat P_K\mathscr Z_t)^{-1}
 \mathscr Z_t^*\widehat E\right]Q_0^{1/2}.
                                                               \tag{NFI10}
\]
Indeed the middle inverse is carried into STF10's inverse by the
unitary $\mathscr R_t$, and its two adjacent factors carry
$E$ into $\widehat E$. This is an equality of matrices on the
original quotient value space, with its original $Q_0$ on both
sides. All cross terms survive the attained minimum. It applies
to the four actual pairs in STF14 and to the target quotient
$(\mathcal U,J)$; the empty quotient conventions remain those
of STF8. Thus it transfers every complement and both allocations,
rather than just their sum of determinants.

Finally use the eight labelled occurrences $(X_a,\sigma_a)$ of
STF16, retaining repeated subspaces separately. The remaining
signed count in STF22 is now the fully original-data expression
\[
 \mathcal B_e^{>R}
  =\sum_{t=q-1,q}\int_1^\infty
       \sum_{a=1}^8\sigma_a
       \#\{\lambda_b(\widehat K_{X_a,t})>u\}
       \,\frac{du}{u}.
                                                               \tag{NFI11}
\]
All eigenvalues count multiplicities. The matrices in NFI11 are
obtained from the finite coefficients NTG4--12, the original
analytic recursion NTG9, and the actual CPQ frames in NFI7.
The equality follows from the proved unitary and partial-isometry
maps, not from a comparison of dimensions. With the accepted prefix
$R_j$, the finite discrepancy bounds STF20--21 and its prefix
budget apply without change. NFI11 gives the full coefficient
formula for the growing-degree quantity; it does not assert a
large-degree estimate for that quantity.
\endgroup
